\section{Problem Summary}
Count the number of binary strings of length \(n\), and report the answer modulo \(10^9 + 7\).

\section{Editorial}
Each position has exactly two choices: \(0\) or \(1\). Therefore the total number of strings is
\[
2^n.
\]

The only real work is computing that power under a modulus. Fast exponentiation is the clean standard tool:
\begin{itemize}
\item repeatedly square the base
\item multiply it into the answer whenever the current bit of the exponent is \(1\)
\end{itemize}

That keeps the arithmetic bounded by the modulus and runs in logarithmic time.

\section{Complexity Analysis}
\begin{itemize}
\item Time: \(O(\log n)\)
\item Memory: \(O(1)\)
\end{itemize}

\section{Notes}
Even though a simple loop of \(n\) doublings would pass here, binary exponentiation scales better and keeps the implementation standard.
